\chapter{Percurso do TCC}

O TCC se desenvolve em etapas. Embora cada curso possa organizar o processo de modo próprio, há um percurso comum: escolha do tema, definição do problema, elaboração do projeto, desenvolvimento da pesquisa, redação, revisão, submissão à banca e entrega final.

\section{Etapas recomendadas}

\begin{longtable}{p{0.23\linewidth}p{0.34\linewidth}p{0.32\linewidth}}
\toprule
Etapa & Produto esperado & Pergunta orientadora \\
\midrule
\endfirsthead
\toprule
Etapa & Produto esperado & Pergunta orientadora \\
\midrule
\endhead
Exploração inicial & Lista de interesses e temas possíveis & O que vale a pena estudar na minha área? \\
Delimitação & Tema recortado e justificativa inicial & O tema é viável no tempo disponível? \\
Problematização & Problema de pesquisa & Que pergunta o trabalho pretende responder? \\
Planejamento & Projeto de pesquisa ou plano de trabalho & Como a pergunta será respondida? \\
Execução & Leituras, coleta, análise ou desenvolvimento & O método está sendo seguido com rigor? \\
Redação & Seções textuais completas e revisadas & O texto comunica o percurso de modo claro? \\
Normalização & Versão conforme padrão institucional & A forma ajuda a leitura e a avaliação? \\
Defesa ou avaliação & Apresentação e arguição, quando houver & O estudante domina o próprio trabalho? \\
Versão final & Arquivo final corrigido & As correções foram incorporadas? \\
\bottomrule
\end{longtable}

\section{Cronograma}

O cronograma deve ser realista. Trabalhos acadêmicos costumam atrasar quando o estudante subestima leitura, coleta de dados, aprovação ética, análise, revisão textual e normalização. É recomendável reservar tempo para imprevistos e para versões sucessivas.

Um cronograma mínimo deve prever:

\begin{itemize}
  \item definição do tema e aceite de orientação;
  \item elaboração ou revisão do projeto;
  \item levantamento bibliográfico;
  \item procedimentos éticos, quando aplicáveis;
  \item coleta, produção ou seleção de materiais;
  \item análise;
  \item redação das seções textuais;
  \item revisão de citações e referências;
  \item revisão formal e normalização;
  \item envio à banca;
  \item ajustes posteriores à avaliação.
\end{itemize}

\caixaatencao{O prazo de entrega à banca não deve ser tratado como prazo de término da escrita. A versão enviada para avaliação precisa chegar madura, revisada e normalizada.}

\section{Registro do processo}

O estudante deve manter registro organizado das decisões tomadas durante o TCC. Isso inclui versões do texto, referências consultadas, critérios de busca, bases utilizadas, instrumentos de coleta, pareceres, autorizações e orientações recebidas.

Esse registro facilita a escrita da metodologia, reduz risco de perda de informação e ajuda a responder perguntas da banca. Também permite reconstruir o percurso caso o trabalho precise ser revisado meses depois.

\section{Relação com a banca}

A banca avalia o trabalho entregue, não a intenção do estudante. Por isso, escolhas importantes precisam estar explícitas no texto. Se o método teve limitação, ela deve ser reconhecida. Se o recorte foi reduzido por viabilidade, isso deve ser justificado. Se determinada fonte foi privilegiada, o critério deve ser apresentado.

Uma boa defesa, quando prevista pelo curso, não compensa uma versão escrita frágil; mas uma apresentação clara pode demonstrar domínio, maturidade e capacidade de argumentação.
